%
% 51-argument.tex
%
% (c) 2026 Prof Dr Andreas Müller
%
\subsection{Argumentänderungen}
Das Argument einer komplexen Zahl $z$ ist der Winkel $\varphi$, für den
$z=|z|e^{i\varphi}$.
Weil die imaginäre Exponentialfunktion wegen $e^{2\pi i}=1$ periodisch
ist, ist das Argument nur bis auf ein ganzzahliges Vielfaches von $2\pi$
festgelegt.
Dies wurde auch schon in
Abschnitt~\ref{buch:analytisch:fortsetzung:subsection:ueberlagerung}
beobachtet, wo die Logarthmusfunktion untersucht wurde.
Der Logarithmus der komplexen Zahl $re^{i\varphi}$ ist
\[
\log (re^{i\varphi}) = \log r + i\varphi + 2\pi i k
\]
für $k\in\mathbb{Z}$.
Die Logarithmusfunktion verspricht also, im Imaginärteil das
Argument zu liefern, versagt aber weil das Resultat nicht eindeutig
bestimmt ist.

%
% Das Argument einer Funktion
%
\subsubsection{Das Argument einer Funktion}
Für $z=0$ ist das Argument undefiniert, es hat bei $0$ eine
Singularität.
Für eine komplexe Funktion $f(z)$ ist es daher ausgeschlossen,
die Funktion $\log f(z)$ zu definieren, wenn die Funktion $f(z)$
eine Nullstelle hat.
Schliesst man Punkte des Definitionsbereichs aus, in denen $f(z)$
verschwindet, dann ist der Wert von $\log f(z)$ bis auf ein ganzzahliges
Vielfaches von $2\pi$ festgelegt.
Es ist also immer noch nicht möglich, einen eindeutigen Wert von $\log f(z)$
festzulegen.
Man kann aber einen der möglichen Werte für $\log f(z_0)$ an der Stelle
$z_0$ zu wählen und dann versuchen, die Werte von $\log f(z)$ für $z$ 
in einer Umgebung von $z_0$ so definieren, dass eine stetige
Funktion entsteht.

\begin{beispiel}
Die Funktion $f(z) = z^3$ bildet $z=re^{i\varphi}$ auf
$f(z) = r^3e^{3i\varphi}$ ab.
Der Logarithmus ist
\[
\log f(z)
=
\log r^3 + 3i\varphi
+
2\pi k
=
3\log r + 3i\varphi
+
2\pi i k
\]
für $k\in\mathbb{Z}$.
Um den Logarithmus festzulegen, muss man also den Wert für einen
Punkt  $z_0$ festzulegen.
Wir wählen $z_0=i=e^{i\frac{\pi}2}$, wir müssen also einen der Werte
\[
\dots
-\frac{9\pi}{2},
-\frac{5\pi}{2},
-\frac{\pi}{2},
{\color{darkred}\frac{3\pi}2}, 
\frac{7\pi}2,
\frac{11\pi}2,
\frac{15\pi}2,
\dots
\]
auswählen.
Der mittlere Wert $3\pi/2$ ist der offensichtliche, aber jede andere
Wahl ist genauso zulässig.

Mit der Wahl $\log f(i) = -\frac{\pi}2i$ können wir jetzt die anderen
Werte des Logarithmus in einer Umgebung von $z_0=i$ berechnen.
Auf dem Gebiet $\mathbb{C}\setminus\{x\in\mathbb{R}\mid x\le 0\}$,
der komplexen Ebene ohne die negative reelle Achse, kann das Argument
von $z=re^{i\varphi}$ im Intervall $\varphi\in(-\pi,\pi)$ eindeutig
festgelegt werden.
Mit dieser Wahl des Argumentes von $z$ ist dann
\[
\log f(z)
=
\log|f(z)| - \frac{\pi}2 + 3\biggl(\varphi-\frac{\pi}2\biggr)i.
\]
Mit dieser Wahl des Wertes von $\log f(z_0)$ haben positive reelle Werte
von $z$ das Argument $\varphi=0$ und damit die Werte
\[
\log f(r) = 3\log r - \frac{\pi}2 -3\frac{\pi}2 = 3\log r-2\pi
\]
des Logarithmus.
Wendet man die Exponentialfunktion an, ergibt sich der reelle Wert
$r^3$, da der Imaginärteil ein ganzzahliges Vielfaches von $2\pi$ ist.
\end{beispiel}

%
% Argument entlang einer Kurve
%
\subsubsection{Argument entlang einer Kurve}
Die Konstruktion von $\log f(z)$ im vorangegangenen Abschnitt
hat das Problem, dass der Logarithmus nicht eindeutig definiert werden kann,
in einer Umgebung eines einzelnen Wertes gelöst.
Verwendet wurde dabei, dass sich das Argument von $\log f(z)$ immer von
einem bekannten Wert $\log f(z_1)$ auf eine Umgebung ausdehnen lässt.
Ausgehend vom Wort in einem anderen Punkt $z_2$ dieser Umgebung, kann
$\log f(z)$ auf eine Umgebung von $z_2$ ausgedehnt werden.
Dies ist genau die Vorgehensweise der analytischen Fortsetzung einer
Funktion, die in Abschnitt~\ref{buch:analytisch:fortsetzung:subsection:kurve}
studiert wurde.
Sie garantiert, dass es entlang einer Kurve $\gamma$ im Definitionsbereich
der Funktion $f(z)$ immer möglich ist, den Logarithmus so zu wählen,
dass $\log f(\gamma(t))$ stetig ist.

Die Festlegung der Werte des Logarithmus von $f(z)$ entlang der Kurve
bedeutet, dass es stetige Funktionen $r(t)$ und $\varphi(t)$ mit
\begin{equation}
\log f(\gamma(t)) = \log r(t) + i\varphi(t)
\label{buch:analytisch:argument:argument:eqn:logf1}
\end{equation}
gibt, die
$f(\gamma(t))
=
e^{r(t)}\bigl(\cos\varphi(t)+i\sin\varphi(t)\bigr)$
erfüllen.

Für eine holomorphe Funktion $f(z)$ und eine differenzierbare Kurve
$\gamma(t)$ sind die Funktionen $r(t)$ und $\varphi(t)$ nicht nur
stetig, sondern sogar differenzierbar.
Man kann daher die Ableitung auf zwei verschiedene Arten berechnen, nämlich
einerseits mit der Kettenregel aus der Zusammensetzung von $\log\circ f$,
und andererseits aus \eqref{buch:analytisch:argument:argument:eqn:logf1}.
Die beiden Arten geben
\[
\frac{d}{dt}\log f(\gamma(t))
=
\begin{cases}
\displaystyle
\frac{f'(\gamma(t))}{f(\gamma(t))}
&\qquad\text{mit der Kettenregel}
\\[10pt]
\displaystyle
\frac{\dot{r}(t)}{r(t)} + i\dot{\varphi}(t)
&\qquad\text{aus 
\eqref{buch:analytisch:argument:argument:eqn:logf1}.}
\end{cases}
\]
Der Realteil beschreibt also die Änderung des Betrags $|f(z)|$ 
während der Imaginärteil die Änderung des Argumentes bestimmt.
Für eine auf dem Intervall $I=[a,b]$ definierte geschlossene Kurve
$\gamma$ kann man beide Formen der Ableitung entlang der Kurve
integrieren, esfolgt dann
\begin{align}
\oint_\gamma \frac{f'(z)}{f(z)}\,dz
&=
\int_I \frac{\dot{r}(t)}{r(t)}\,dt
+
i\int_I \dot{\varphi}(t)\,dt
\notag
\\
&=
\Bigl[ \log r(t) \Bigr]_a^b
+
i\Bigl[\varphi(t)\Bigr]_a^b
\notag
\\
&=
\bigl(\log r(b) - \log r(a)\bigr)
+
i
\bigl(\varphi(b)-\varphi(a)\bigr).
\label{buch:analytisch:argument:eqn:argumentintegral0}
\end{align}
Da die Kurve geschlossen ist, ist $r(b)=r(a)$, der Realteil von
\eqref{buch:analytisch:argument:eqn:argumentintegral0}
verschwindet daher.

\begin{satz}[Argumentintegral]
Für eine holomorphe Funktioin $f(z)$ und eine differenzierbare Kurve
$\gamma(t)$ im Definitionsgebiet von $f(z)$, auf der $f(z)$ nicht verschwindet,
ist das Integral
\begin{equation}
\oint_\gamma \frac{f'(z)}{f(z)}\,dz
=
2\pi i k
\label{buch:analytisch:argument:eqn:argumentintegral}
\end{equation}
mit $k\in\mathbb{Z}$.
\index{Argumentintegral}%
\end{satz}

Wir nennen \eqref{buch:analytisch:argument:eqn:argumentintegral}
das Argumentintegral.
Es hängt offenbar nicht von der konkreten Wahl des Logarithmus ab.
Insbesondere kann es auch ohne die Aufwendige konstruktion einer
analytischen Fortsetzung der Funktion $\log f(z)$ berechnet werden.


